% ============================================================================
% Fichier  : front/01_avant_propos.tex
% Titre    : Avant-propos
% Auteur   : MINKA MI NGUIDJOÏ Thierry Emmanuel
% Version  : 0.1 -- novembre 2026
% ============================================================================
% 8 jours de 4 h (32 h au programme, 30 h effectives : les 2 dernières heures
% du J08 sont offertes). Seuls A0 et A1 (HAL) sont cités comme travaux de
% l'auteur.
% ============================================================================

\chapter*{Avant-propos}
\addcontentsline{toc}{chapter}{Avant-propos}
\markboth{Avant-propos}{Avant-propos}

\epigraph{« Le contrôle interne ne disparaît pas quand l'administration
s'informatise~; il change de place, et l'auditeur doit le suivre. »}%
{\textit{Principe du cours}}

Ce cours s'adresse aux auditeurs et aux contrôleurs des finances publiques
du Programme Supérieur de Spécialisation en Finances Publiques. Ils
connaissent déjà le contrôle sur pièces~: le mandat, la dépense, la
régularité, le contradictoire. Ils vont désormais exercer devant des
applications de budget, de dépense, de paie et de recette, où la pièce est
un enregistrement, la signature un droit d'accès, et le visa un contrôle
automatique.

Il repose sur une conviction~: l'informatique ne supprime pas le contrôle
interne, elle le déplace. La séparation entre l'ordonnateur et le comptable
devient une matrice de droits~; la vérification du service fait devient un
contrôle de saisie~; l'archive devient une sauvegarde. Un auditeur qui
continue de ne regarder que le dossier papier ne voit plus qu'une partie du
système. Le cœur du cours est donc le \textbf{système de contrôle interne}
et son \textbf{évaluation}~: le décrire, le tester, classer ses
déficiences, conclure sur son efficacité.

\bigskip

\textbf{Le cadre camerounais.} Au Cameroun, dans le secteur public, c'est
le Contrôle Supérieur de l'État qui a la compétence explicite du contrôle
des systèmes d'information~: son texte organique, le décret n°~2013/287
du 4~septembre~2013, la range parmi les cinq contrôles qu'il effectue.
L'ANTIC conduit les audits de sécurité des réseaux et systèmes
d'information au titre de la loi n°~2010/012. La Chambre des Comptes, juge
des comptes publics, peut s'y retrouver par extension, à travers les
comptes que ces systèmes produisent. Le Chapitre~\ref{chap:institutions}
précise ces compétences.

\bigskip

\textbf{Ce que ce cours n'est pas.} Il n'est pas un cours d'informatique,
ni une préparation à une certification, ni un cours de sécurité offensive.
Il ne demande aucun prérequis technique~: un tableur suffit pour les labs.

\bigskip

\textbf{Ce que ce cours est.} Un traité de quinze chapitres, en six
parties précédées d'un prologue, conçu pour trente heures effectives en
huit journées de quatre heures. La théorie y est lisible~; chaque notion
sert un geste de mission. La Partie~I pose le contrôle interne et la
manière de l'évaluer. La Partie~II présente les institutions et les
normes. La Partie~III conduit la mission, de l'ordre de mission au cycle
suivant~: principes, décomposition, instrument \OUI{}~/ \NON{}~/ \NEK{},
mesures de la conformité et de l'incertitude, rapport et suivi. La
Partie~IV évalue le contrôle interne domaine par domaine~: accès,
exploitation, applications, données, projets. La Partie~V réunit tout le
cours dans un cas intégrateur traité en salle. La Partie~VI ouvre sur la
recherche et propose des sujets de mémoire.

\bigskip

\textbf{Une méthode fondée sur une recherche.} La démarche de mission est
celle du cours d'audit des systèmes d'information de l'auteur à l'École
Nationale Supérieure Polytechnique de Yaoundé, elle-même instanciation,
pour le praticien, du programme de recherche de
l'\emph{auditologie}~\cite{minka2026auditologie}. L'instrument de réponse
\OUI{}~/ \NON{}~/ \NEK{} et ses mesures viennent de l'article
\cite{minka2026nek}. Ces deux textes sont des prépublications déposées sur
HAL, et le cours les cite comme telles.

\bigskip

\textbf{Trois entités fictives.} Le cours suit trois dossiers. Au
\emph{MINÉQUIR}, ministère imaginaire de l'Équipement rural, une facture
est payée deux fois par l'application de dépense. À la \emph{CAPNO},
caisse publique de prévoyance de Garoua, le fichier de paie contient des
agents qui n'existent pas. À la recette des finances de Bafoussam, le
projet de dématérialisation \emph{e-RECETTES} a été livré sans que
personne n'en ait vérifié les sauvegardes. Les noms, les dates et les
montants sont fictifs~; les mécanismes sont exacts.

\bigskip

\textbf{Comment travailler.} Entre deux journées, faites la fiche
d'auto-évaluation de chaque chapitre et les labs de NDAP SUKULU qui s'y
rattachent. Prenez surtout l'habitude que le cours veut transmettre~:
séparer ce que l'on a constaté de ce que l'on en pense, et écrire ce que
l'on ne sait pas.

\bigskip

\begin{flushright}
    \IfFileExists{../../Logos/signature_avant_propos.png}{\includegraphics[height=1.6cm]{../../Logos/signature_avant_propos.png}}{\textit{MINKA MI NGUIDJOÏ Thierry Emmanuel}}\\
    \textit{Yaoundé, novembre 2026}
\end{flushright}
